\documentclass[12pt,a4paper]{article}

% Packages
\usepackage[utf8]{inputenc}
\usepackage[T1]{fontenc}
\usepackage{amsmath}
\usepackage{amssymb}
\usepackage{hyperref}
\usepackage{graphicx}
\usepackage{booktabs}

\title{A Sample \LaTeX{} Document}
\author{Jane Doe \and John Smith}
\date{\today}

\begin{document}

\maketitle

% Table of contents
\tableofcontents
\newpage

%% Introduction
\section{Introduction}
\label{sec:intro}

This document demonstrates the syntax of \LaTeX{}, a document preparation system
based on Donald Knuth's \TeX{} typesetting engine. It is widely used in academia
for \textbf{scientific}, \textit{mathematical}, and \emph{technical} writing.

See Section~\ref{sec:math} for examples of mathematical typesetting.

\subsection{Background}
\label{sec:background}

\LaTeX{} was developed by Leslie Lamport in 1983 as a higher-level interface to
\TeX{}. It provides macros for document structuring, cross-referencing, and
bibliography management~\cite{lamport1994latex}.

\subsubsection{Why Use \LaTeX{}?}
\begin{itemize}
  \item Superior mathematical typesetting
  \item Automatic numbering and cross-references
  \item Consistent, professional output
  \item Free and open source
\end{itemize}

\section{Mathematics}
\label{sec:math}

\subsection{Inline and Display Math}

Inline math uses dollar signs: $E = mc^2$ and $\sum_{i=1}^{n} i = \frac{n(n+1)}{2}$.

Display math uses the \texttt{equation} environment:

\begin{equation}
  \label{eq:euler}
  e^{i\pi} + 1 = 0
\end{equation}

Unnumbered display equations use \verb|\[...\]|:

\[
  \int_{-\infty}^{\infty} e^{-x^2} \, dx = \sqrt{\pi}
\]

\subsection{Aligned Equations}

\begin{align}
  f(x) &= x^2 + 2x + 1 \\
       &= (x + 1)^2
\end{align}

\subsection{Theorem Environments}

\newtheorem{theorem}{Theorem}
\newtheorem{lemma}[theorem]{Lemma}

\begin{theorem}[Pythagorean Theorem]
  For a right triangle with legs $a$ and $b$ and hypotenuse $c$:
  \[
    a^2 + b^2 = c^2
  \]
\end{theorem}

\section{Environments}
\label{sec:environments}

\subsection{Tables}

\begin{table}[ht]
  \centering
  \caption{Comparison of typesetting systems}
  \label{tab:comparison}
  \begin{tabular}{lcc}
    \toprule
    System    & Math support & Free \\
    \midrule
    \LaTeX{}  & Excellent    & Yes  \\
    Typst     & Good         & Yes  \\
    Word      & Limited      & No   \\
    \bottomrule
  \end{tabular}
\end{table}

\subsection{Figures}

\begin{figure}[ht]
  \centering
  % \includegraphics[width=0.5\textwidth]{example-image}
  \caption{An example figure.}
  \label{fig:example}
\end{figure}

\subsection{Verbatim}

\begin{verbatim}
#include <stdio.h>
int main() {
    printf("Hello, World!\n");
    return 0;
}
\end{verbatim}

\subsection{Code Listings}

\begin{lstlisting}[language=Python]
def greet(name):
    print(f"Hello, {name}!")
\end{lstlisting}

\subsection{Math Sub and Superscripts}

The quadratic formula: $x_{1,2} = \frac{-b \pm \sqrt{b^2 - 4ac}}{2a}$

Einstein's mass-energy: $E = mc^2$

\subsection{TODO Notes}

\todo{Add more examples here.}

\section{Cross-References and URLs}
\label{sec:refs}

Refer to Equation~\ref{eq:euler}, Table~\ref{tab:comparison},
and Figure~\ref{fig:example}.

Visit \url{https://www.latex-project.org} or use
\href{https://ctan.org}{CTAN} for packages.

% Define a custom command
\newcommand{\R}{\mathbb{R}}
\newcommand{\norm}[1]{\left\|#1\right\|}

The set of real numbers is $\R$ and a norm is $\norm{x}$.

% Bibliography
\bibliographystyle{plain}
\begin{thebibliography}{9}
  \bibitem{lamport1994latex}
    Leslie Lamport,
    \textit{\LaTeX: A Document Preparation System},
    Addison-Wesley, 1994.

  \bibitem{knuth1984tex}
    Donald E. Knuth,
    \textit{The \TeX book},
    Addison-Wesley, 1984.
\end{thebibliography}

\end{document}
